\documentclass[10pt,a4paper]{amsart}
\usepackage[margin=19mm]{geometry}
\usepackage{amsmath,amssymb,mathtools}
\usepackage[scr=rsfs,cal=euler]{mathalfa}
\usepackage{xfrac}
\usepackage[unicode,hidelinks,bookmarks=false]{hyperref}
\newcommand{\ie}{i.e.\ }
\newcommand{\cf}{cf.\ }
\newcommand{\resp}{respectively\ }
\newcommand{\Subsection}{\subsection}
\providecommand{\nobreakdash}{\nobreak\hbox{-}}
\setlength{\emergencystretch}{2em}
\multlinegap0pt
\usepackage{bm}
\usepackage{amssymb}
\usepackage{mathtools}
\usepackage{cancel}
\usepackage{xfrac}
\usepackage[shortlabels]{enumitem}
\usepackage{xcolor}
\usepackage{tikz}
\usepackage{float}
\newtheorem{theorem}{Theorem}
\newtheorem{lemma}{Lemma}
\newtheorem{corollary}{Corollary}
\title{Global synchronization beyond dense graphs: the case of threshold graphs}
\author{Hongjin Wu, Ulrik Brandes}
\date{Source: arXiv:2511.12646v6 (CC BY 4.0)}
\begin{document}
\maketitle

\noindent\emph{Source and licence.} Global synchronization beyond dense graphs: the case of threshold graphs. Version arXiv:2511.12646v6 (CC BY 4.0), distributed by arXiv under the Creative Commons Attribution 4.0 International licence, http://creativecommons.org/licenses/by/4.0/, as recorded in the arXiv licence metadata for this exact version and shown on the abstract page https://arxiv.org/abs/2511.12646v6. Full licence notice: \texttt{licenses/arxiv-2511.12646-CC-BY-4.0.txt}.\par\smallskip

\noindent\textit{Editorial scope.} Counted unit: the article's theorem main (source0.tex lines 525--527). The article's complete original proof of it is delivered with it (source0.tex lines 1814--1862, one \texttt{proof} environment of 49 lines, which the article places away from the statement). No mathematics is written, repaired or re-derived: the statement and the proof are the article's own lines, in the article's own notation, and no retained line is substituted or re-flowed.  Retained uncounted as the source-local context the proof needs: lines 417--421; lines 1015--1018; lines 1180--1181; lines 1222--1225; lines 1400--1425.  Nothing was omitted from any retained span, so the package carries no omission record.  No retained line contains a citation, so the wrapper carries no bibliography and no bibliography entry.  No retained line contains a figure environment, an image-inclusion command or drawing code, so no image is delivered and the package declares no asset dependency.  Editorial formatting only, in addition: an AMS \texttt{amsart} preamble that restates the article's own declarations and macros used by the retained lines, namely the environments \texttt{theorem}, \texttt{lemma}, \texttt{corollary} on separate counters, together with the standard packages those lines need.  The retained environments print in this document as Theorem 1, Lemma 1, Corollary 1 in order of appearance, whereas the article prints the same items with the numbering of its own full document.  The complete native source of the article as distributed in the arXiv e-print for this exact version is bundled unchanged as \texttt{sources/189215/source0.tex}.
\begin{equation}\label{E}
\begin{aligned}
    E_G(\boldsymbol\theta) := &\frac{1}{2} \sum_{1\leq i,j\leq n} \boldsymbol{A}_{ij} \big(1-\cos(\theta_i - \theta_j)\big). \\
    \end{aligned}
\end{equation}

\begin{lemma}\label{stable-eq}
If $\boldsymbol{\theta}\in\mathbb R^n$ is a second-order stationary point of~\eqref{E},
then for every $i\in V$ there exists $\mu_i\ge 0$ such that $$\sum_{j\in N(i)} \boldsymbol{v}_j = \mu_i \boldsymbol{v}_i.$$
\end{lemma}

\begin{equation}\label{st-closed}
N[i] = N[j].\end{equation}

\begin{corollary}\label{cor:twins-sync}
Let $\boldsymbol{\theta}\in\mathbb R^n$ be a second-order stationary point of the energy function $E_G(\boldsymbol{\theta})$. 
If nodes $a$ and $b$ form a pair of closed twins (satisfying \eqref{st-closed}), then $\boldsymbol{v}_a = \boldsymbol{v}_b$.
\end{corollary}

\begin{lemma}\label{lemma:pendant}
Let $G=(V,E)$ be a graph, and let $W \subseteq V$ admit a partition
\[ W = Q \sqcup S_1 \sqcup S_2 \sqcup P\]
with the following structural conditions:

(i) Every node in $Q$ has all its neighbors inside $S_1$, 
    i.e., $N(i) \subseteq S_1$ for all $i \in Q$;

(ii) For each $i \in S_1 \sqcup S_2$, its neighborhood is 
    $P \sqcup N_Q(i) \sqcup (S \setminus \{i\})$, 
    where $$N_Q(i) = \{\, q \in Q : q \sim i \,\}.$$
    
(iii) For any \( i \in S_1 \sqcup S_2 \), its neighborhood can be decomposed into three parts: the set \(P\), its neighbours inside \(Q\), and the rest of the clique \(S_1 \sqcup S_2\).  
That is
\[
N(i) \;=\; P \sqcup N_Q(i) \sqcup \bigl((S_1 \sqcup S_2)\setminus \{i\}\bigr).
\]
Assume further that \( \{\boldsymbol{v}_i\}_{i\in V} \) corresponds to a second-order stationary point
\( \boldsymbol{\theta} \) of the Kuramoto energy~\eqref{E}, at which the induced subgraph on
\( Q \sqcup S_1 \) is synchronized, i.e.,
\[
\boldsymbol{v}_i = \boldsymbol{v}, \quad \forall i \in Q \sqcup S_1 ,
\]
for some common phasor \( \boldsymbol{v} \).
Then the set $S_1\sqcup S_2$ synchronize.
\end{lemma}

\begin{theorem}\label{main}
Every connected threshold graph is second-order globally synchronizing. In particular, every connected threshold graph is globally synchronizing.
\end{theorem}

\begin{proof}[Proof of Theorem~\ref{main}]
Any connected threshold graph can be described by a bit sequence of $1$'s and $0$'s, where the last bit must be a $1$; otherwise, the threshold graph is not connected.
For simplicity, we decompose the sequence into maximal contiguous blocks of $1$'s and $0$'s:
\[
U_1,\, I_2,\, U_2,\, \dots,\, I_{k-1},\, U_{k-1},\, I_k,\, U_k,
\]
or
\[
I_1,\, U_1,\, I_2,\, U_2,\, \dots,\, I_{k-1},\, U_{k-1},\, I_k,\, U_k,
\]
where $U_i$ denotes a subsequence of $1$'s and $I_i$ denotes a subsequence of $0$'s.
Without loss of generality, we focus on threshold graphs $G=(V,E)$ of the first sequence type; the second case follows by an entirely analogous argument.

Let $\boldsymbol{\theta}$ be a second-order stationary point of the energy function $E_G(\boldsymbol{\theta})$ of the homogeneous Kuramoto model on $G$.
We prove that $\boldsymbol{\theta}$ is a fully synchronized state by induction on $m$, where $m$ denotes the number of $1$-blocks counted from the end of the threshold sequence.

For the base case, observe that all nodes within a block of $1$'s form a set of closed twins.
Hence, at any second-order stationary point $\boldsymbol{\theta}$, the nodes in the last block $U_k$ must synchronize, as follows directly from Corollary~\ref{cor:twins-sync}.

Assume as the inductive hypothesis that the nodes in the last $m$ blocks of $1$'s, namely $U_{k-m+1}, \ldots, U_k$, are synchronized, i.e., there exists a common phasor $\boldsymbol{v}$ such that $\boldsymbol{v}_i = \boldsymbol{v}$ for all $i \in U_{k-m+1} \sqcup \cdots \sqcup U_k$. Lemma~\ref{stable-eq} then implies that the nodes in the adjacent $0$-blocks $I_{k-m+1} \sqcup \cdots \sqcup I_k$ also satisfy $\boldsymbol{v}_i = \boldsymbol{v}$.

Define
\[
\begin{aligned}
Q   &:= I_{k-m+1} \sqcup \cdots \sqcup I_k, \\
S_1 &:= U_{k-m+1} \sqcup \cdots \sqcup U_k, \\
S_2 &:= U_{k-m}, \\
P   &:= U_1 \sqcup \cdots \sqcup U_{k-m-1}
         \sqcup I_2 \sqcup \cdots \sqcup I_{k-m}.
\end{aligned}
\]
We now verify that the assumptions of Lemma~\ref{lemma:pendant} are satisfied by the above
decomposition, which follows directly from the threshold construction.

By construction, the vertices in $Q$ correspond to the most recent $0$-blocks
preceding $S_1$. Hence, every vertex in $Q$ is adjacent only to vertices in $S_1$.
Moreover, since $S_1$ and $S_2$ consist of consecutive dominating blocks,
the induced subgraph on $S_1 \cup S_2$ is a clique.
Finally, for each $i\in S_1$, its neighborhood can be decomposed as
\[
N(i) = P \;\sqcup\; \bigl((S_1\cup S_2)\setminus\{i\}\bigr) \;\sqcup\; N_Q(i),
\]
where $P$ denotes the set of vertices added earlier in the construction.
Hence, all the assumptions of Lemma~\ref{lemma:pendant} are satisfied,
and we conclude that the sets $S_1$ and $S_2$ synchronize.
In other words, the last $m+1$ blocks of $1$'s $U_{k-m}, U_{k-m+1}, \cdots, U_k$ synchronize.

Therefore, every second-order stationary point is a synchronous state, and it follows that $G$ is second-order globally synchronizing, and in particular, it is globally synchronizing.
\end{proof}

\end{document}
